\chapter[Separability, Normality, and Finite Fields]{Separability, Normality, and Finite Fields}
\chaptermark{Separability and Finite Fields}

Embeddings count genuinely different conjugate roots.  Separability measures
that distinction; normality says that no conjugate root escapes the field.
Only after both ideas are clear will we define a Galois extension.

\section{Separable Polynomials and Extensions}
\begin{definition}
A nonconstant polynomial \(f\in F[x]\) is \emph{separable} if its roots in a
splitting field are distinct.  An algebraic element is \emph{separable} over
\(F\) when its minimal polynomial is separable; an algebraic extension is
separable when each of its elements is separable.
\end{definition}

The definition applies to arbitrary polynomials.  For an irreducible
polynomial it asks whether the possible conjugates of one root occur only
once.  This is the distinction needed for counting embeddings.

\begin{proposition}[Derivative criterion]\label{prop:separable-criterion}
For nonconstant \(f\in F[x]\),
\[
f\text{ has a repeated root in }\overline F
\quad\Longleftrightarrow\quad
\gcd(f,f')\ne1.
\]
For irreducible \(f\), this is equivalent to \(f'=0\).
\end{proposition}
\begin{proof}
Write \(f=(x-a)^rg\) in a splitting field, with \(g(a)\ne0\).  The product
rule gives
\[
f'=(x-a)^{r-1}\bigl(rg+(x-a)g'\bigr).
\]
Thus \(a\) is repeated exactly when \(r\ge2\), equivalently when
\(f(a)=f'(a)=0\).  A common divisor of \(f\) and \(f'\) is therefore
equivalent, after splitting, to a common linear factor, hence to a repeated
root.  For irreducible \(f\), the gcd is either \(1\) or an associate of
\(f\); the latter can occur only if \(f'=0\).
\end{proof}

\begin{remark}
In characteristic \(p>0\), \(f'=0\) exactly when every exponent appearing
in \(f\) is divisible by \(p\), so
\[
f(x)=g(x^p).
\]
This is the source of inseparability.  For example,
\[
x^p-t\in\mathbb F_p(t)[x]
\]
has derivative zero and is irreducible because a rational function whose
\(p\)-th power is \(t\) would have every zero and pole order divisible by
\(p\), whereas \(t\) has a simple zero.  In an algebraic closure it is
\[
x^p-t=(x-t^{1/p})^p;
\]
it has one distinct root, with multiplicity \(p\).
\end{remark}

\begin{example}[A separable polynomial with repeated-looking coefficients]\label{ex:sep-cubic}
Over \(\mathbb F_2\), the polynomial
\[
x^3+x+1
\]
has derivative \(x^2+1\).  It has no root in \(\mathbb F_2\), hence is
irreducible, and it cannot share a root with its derivative.  It is therefore
separable.  The contrast with \(x^2-t\) over \(\mathbb F_2(t)\), whose
derivative is zero, shows that characteristic alone is not the issue:
inseparability comes from the special \(p\)-th-power shape.
\end{example}

\section{Embeddings, Perfect Fields, and Primitive Elements}

The next three ideas answer complementary questions about a finite extension.
Separability asks whether the degree is fully visible through distinct
conjugates; normality, introduced later, asks whether those conjugates remain
inside the field; and the primitive element theorem says that a finite
separable extension can often be described by one well-chosen generator.
The comparison is useful to keep in view:
\[
\begin{array}{c|c}
\text{property} & \text{what it controls}\\ \hline
\text{separable} & \text{distinct embeddings and conjugate roots}\\
\text{normal} & \text{whether all conjugates stay in the extension}\\
\text{Galois} & \text{both conditions, hence internal symmetries}
\end{array}
\]

\begin{definition}
A field \(F\) is \emph{perfect} if every algebraic extension of \(F\) is
separable, equivalently if every irreducible polynomial in \(F[x]\) is
separable.
\end{definition}
\begin{proposition}
Every characteristic-zero field is perfect.  A field \(F\) of characteristic
\(p>0\) is perfect if and only if Frobenius
\[
F\longrightarrow F,\qquad a\longmapsto a^p
\]
is surjective.  In particular every finite field is perfect.
\end{proposition}
\begin{proof}
In characteristic zero a nonzero irreducible polynomial cannot have zero
derivative, so the criterion proves separability.  In characteristic \(p\),
an inseparable irreducible polynomial has derivative zero and is of the form
\(g(x^p)\).  Surjectivity of Frobenius makes each coefficient a \(p\)-th
power, so this is a \(p\)-th power polynomial, impossible for a nonconstant
irreducible polynomial.  Conversely, if \(a\) has no \(p\)-th root, then
\(x^p-a\) has derivative zero and an irreducible factor which is inseparable.
Frobenius is injective on every field and hence bijective on a finite field.
\end{proof}

Before turning this count into a definition, notice the strategy: build the
field one algebraic generator at a time.  At each step an embedding has only
as many possible extensions as the minimal polynomial has roots.  The Tower
Law then turns this local bound into a global one.

\begin{theorem}[Embedding-count inequality]\label{thm:embedding-count}
For finite \(E/F\),
\[
\#\operatorname{Hom}_F(E,\overline F)\le [E:F],
\]
\end{theorem}
\begin{proof}
Choose generators
\[
F=F_0\subseteq F_1\subseteq\cdots\subseteq F_r=E,
\qquad
F_i=F_{i-1}(\alpha_i).
\]
An embedding \(F_{i-1}\to\overline F\) extends across \(\alpha_i\) only by
sending it to a root of the transported minimal polynomial.  By the
root-extension lemma, there are at most
\([F_i:F_{i-1}]\) choices.  Successively extending embeddings therefore
gives
\[
\#\operatorname{Hom}_F(E,\overline F)
\leq\prod_i[F_i:F_{i-1}]=[E:F].
\]
Equality at every step holds exactly when each transported minimal polynomial
has its full number of distinct roots.  We analyze that equality after
introducing separable degree.
\end{proof}

\begin{definition}
For a finite extension, the number
\[
[E:F]_s=\#\operatorname{Hom}_F(E,\overline F)
\]
is its \emph{separable degree}.  Thus
\[
[E:F]_s\leq[E:F],
\]
by Theorem~\ref{thm:embedding-count}.
\end{definition}

The ordinary degree is multiplicative in a tower by the Tower Law.  The
following result says that the count of genuinely distinct embeddings has the
same multiplicative behavior.  Its proof is a refined version of the
preceding counting argument: first choose an embedding of the middle field,
then count its possible continuations.

\begin{proposition}[Multiplicativity of separable degree]\label{prop:separable-degree}
If
\[
F\subseteq K\subseteq E
\]
is a tower of finite extensions, then
\[
[E:F]_s=[E:K]_s[K:F]_s.
\]
\end{proposition}
\begin{proof}
Choose a tower of simple extensions
\[
K=K_0\subseteq K_1\subseteq\cdots\subseteq K_r=E,
\qquad K_i=K_{i-1}(\alpha_i),
\]
and let \(d_i\) be the number of distinct roots of the minimal polynomial
of \(\alpha_i\) over \(K_{i-1}\).  The root-extension lemma shows that an
embedding of \(K_{i-1}\) has exactly \(d_i\) extensions to \(K_i\).
Indeed, the transported minimal polynomial is obtained by applying the
embedding to its coefficients.  It has the same degree and its derivative
is the transported derivative, so its gcd with its derivative has the same
degree; hence it has exactly the same number \(d_i\) of distinct roots.

It follows that every \(F\)-embedding \(\sigma:K\to\overline F\) has
\[
\prod_{i=1}^r d_i=[E:K]_s
\]
extensions to \(E\).  The equality with \([E:K]_s\) is obtained by applying
the same step-by-step count beginning with the identity embedding of \(K\).
Counting first the possible restrictions to \(K\), then their extensions,
proves the formula.
\end{proof}

\begin{theorem}[Separable degree criterion]\label{thm:separable-degree}
For a finite extension \(E/F\),
\[
[E:F]_s=[E:F]
\]
if and only if \(E/F\) is separable.
\end{theorem}
\begin{proof}
Choose a simple generating tower as in the proof of
Proposition~\ref{prop:separable-degree}.  The embedding count there gives
\[
[E:F]_s=\prod_i d_i,
\qquad
[E:F]=\prod_i[K_i:K_{i-1}].
\]
If \(E/F\) is separable, the minimal polynomial of \(\alpha_i\) over
\(K_{i-1}\) divides its separable minimal polynomial over \(F\), so it is
separable.  Thus \(d_i=[K_i:K_{i-1}]\) for every \(i\), proving equality.

Conversely, equality of the two products forces equality in every factor,
so each \(\alpha_i\) is separable over \(K_{i-1}\).  To see that every
\(\gamma\in E\) is separable over \(F\), set \(L=F(\gamma)\).  By
Proposition~\ref{prop:separable-degree},
\[
[E:F]_s=[E:L]_s[L:F]_s.
\]
Both factors are bounded above by their ordinary degrees, while their
product is \([E:F]=[E:L][L:F]\).  Hence
\[
[L:F]_s=[L:F].
\]
For the simple extension \(L=F(\gamma)\), this says precisely that the
minimal polynomial of \(\gamma\) has its full number of distinct roots.
Thus \(\gamma\) is separable.
\end{proof}

\begin{proposition}[Separability in towers]\label{prop:separable-towers}
For a finite tower
\[
F\subseteq K\subseteq E,
\]
\(E/F\) is separable if and only if both \(E/K\) and \(K/F\) are separable.
The same statement holds for arbitrary algebraic towers.
\end{proposition}
\begin{proof}
If \(E/F\) is separable, elements of \(K\) are already separable over
\(F\), and a minimal polynomial over \(K\) divides the corresponding
separable minimal polynomial over \(F\); hence both subextensions are
separable.  Conversely, if both are separable, the multiplicativity formula
and Theorem~\ref{thm:separable-degree} give
\[
[E:F]_s=[E:K][K:F]=[E:F].
\]
Thus \(E/F\) is separable.  For arbitrary algebraic towers, apply the finite
statement as follows.  The forward direction is immediate for \(K/F\), since
\(K\subseteq E\); and, for \(\alpha\in E\), the minimal polynomial over
\(K\) divides the separable minimal polynomial over \(F\), proving that
\(E/K\) is separable.

Conversely, suppose \(K/F\) and \(E/K\) are separable, and choose
\(\alpha\in E\).  The minimal polynomial
\[
m_{\alpha,K}(x)
\]
has only finitely many coefficients.  Let \(K_0\) be the subfield of \(K\)
generated over \(F\) by those coefficients.  Then \(K_0/F\) is finite; it is
separable because every element of the subfield \(K_0\subseteq K\) is already
separable over \(F\).  The polynomial \(m_{\alpha,K}\) lies in \(K_0[x]\) and is
separable because \(E/K\) is separable.  Hence the minimal polynomial of
\(\alpha\) over \(K_0\), which divides \(m_{\alpha,K}\), is separable.
Thus \(K_0(\alpha)/K_0\) is finite separable.  The finite tower statement,
applied to
\[
F\subseteq K_0\subseteq K_0(\alpha),
\]
shows that \(\alpha\) is separable over \(F\).  Since \(\alpha\) was
arbitrary, \(E/F\) is separable.
\end{proof}

\begin{lemma}[Cyclicity lemma]\label{lem:finite-mult-cyclic}
Every finite subgroup of the multiplicative group of a field is cyclic.
\end{lemma}
\begin{proof}
Let \(G\) be such a subgroup and let \(m\) be its exponent.  Every element
of \(G\) is a root of \(x^m-1\), so \(|G|\leq m\).  Conversely, by the
structure theorem for finite abelian groups, \(G\) has an element of order
equal to its exponent \(m\), hence \(m\leq|G|\).  Therefore \(m=|G|\), and
an element of order \(m\) generates \(G\).
\end{proof}

The primitive element theorem is not merely a convenient notation theorem.
It reduces a finite separable extension to a single minimal polynomial, so
that its embeddings and conjugates can be studied through one element.  Over
an infinite field, the proof chooses a linear combination
\(\alpha+c\beta\) avoiding the finitely many values of \(c\) that would make
two embeddings coincide; over a finite field, cyclicity of the multiplicative
group supplies a generator instead.

\begin{theorem}[Primitive element theorem]
Every finite separable extension is simple.
\end{theorem}
\begin{proof}
First suppose \(F\) is infinite and \(E=F(\alpha,\beta)\).  List the
\([E:F]\) \(F\)-embeddings of \(E\) into an algebraic closure as
\(\sigma_i\), which exist by Theorem~\ref{thm:separable-degree}.  For
each pair of distinct images, the equality
\[
\sigma_i(\alpha)+c\sigma_i(\beta)
=\sigma_j(\alpha)+c\sigma_j(\beta)
\]
excludes at most one value of \(c\in F\).  Choose \(c\) outside this finite
set and put \(\gamma=\alpha+c\beta\).  The conjugates of \(\gamma\) are then
\([E:F]\) distinct values.  They are values of distinct restrictions of the
\(\sigma_i\) to \(F(\gamma)\): if two restrictions agreed, they would give
the same displayed value.  Hence
\[
\#\operatorname{Hom}_F(F(\gamma),\overline F)\ge [E:F].
\]
The embedding-count theorem gives the reverse bound by
\([F(\gamma):F]\), while the Tower Law gives
\([F(\gamma):F]\le [E:F]\).  Thus equality holds throughout.
Since \(F(\gamma)\subseteq E\), the Tower Law gives \(E=F(\gamma)\).
Induction reduces finitely many generators to this case.

If \(F\) is finite, then \(E^\times\) is a finite subgroup of a field and
is cyclic by Lemma~\ref{lem:finite-mult-cyclic}.  A generator \(\gamma\) of
\(E^\times\) has \(E=F(\gamma)\).
\end{proof}

\begin{example}[A primitive element]
The extension
\[
\mathbb Q(\sqrt2,\sqrt3)/\mathbb Q
\]
is finite and separable.  The proof above predicts a primitive element; here
one may take \(\gamma=\sqrt2+\sqrt3\).  Since
\[
\gamma^2=5+2\sqrt6,
\]
we recover \(\sqrt6\), and then
\[
\sqrt2=\frac{\gamma^2-1}{2\gamma},
\qquad
\sqrt3=\gamma-\sqrt2.
\]
Thus \(\mathbb Q(\sqrt2,\sqrt3)=\mathbb Q(\sqrt2+\sqrt3)\).
\end{example}

\section{Normal Extensions and Galois Extensions}
\begin{definition}
An algebraic extension \(E/F\) is \emph{normal} if every irreducible
polynomial in \(F[x]\) having one root in \(E\) splits over \(E\).  It is
\emph{Galois} if it is algebraic, normal, and separable.
\end{definition}

The guiding slogan is:
\[
\text{normal means that no conjugate root escapes.}
\]
Thus normality is a condition on the \emph{whole} field, whereas separability
is a condition on the distinctness of roots.  The normality criteria below
translate this root-theoretic statement into the embedding language developed
above.

\begin{theorem}[Normality criteria]\label{thm:normality-criteria}
For finite \(E/F\), the following are equivalent:
\begin{enumerate}
\item \(E/F\) is normal;
\item every \(F\)-embedding \(E\to\overline F\) has image \(E\);
\item \(E\) is a splitting field over \(F\) of a collection of polynomials.
\end{enumerate}
\end{theorem}
\begin{proof}
Assume first that \(E/F\) is normal.  If \(\sigma:E\to\overline F\) is an
\(F\)-embedding and \(\alpha\in E\), then \(\sigma(\alpha)\) is a conjugate
root of \(m_{\alpha,F}\).  Normality makes this root belong to \(E\); hence
\(\sigma(E)\subseteq E\).  Both fields have degree \([E:F]\) over \(F\), so
\(\sigma(E)=E\).

Conversely, suppose every embedding has image \(E\).  If an irreducible
\(f\in F[x]\) has a root \(\alpha\in E\), every root \(\beta\) of \(f\) gives
an embedding \(F(\alpha)\to\overline F\) by Chapter~10.  Extend it to \(E\)
by Theorem~\ref{thm:embedding-extension}.  Its image is \(E\), so
\(\beta\in E\).  Thus \(f\) splits in \(E\), and \(E/F\) is normal.

Finally, suppose \(E\) is a splitting field over \(F\).  Every
\(F\)-embedding of \(E\) fixes the defining coefficients and permutes the
corresponding roots.  Its image is consequently contained in \(E\); equality
follows from degrees, and the already proved implication shows that \(E/F\)
is normal.  If \(E/F\) is
finite normal, choose generators \(\alpha_1,\ldots,\alpha_r\).  The splitting
field in \(E\) of the finitely many minimal polynomials
\(m_{\alpha_i,F}\) contains every generator and is therefore \(E\).
\end{proof}

\begin{example}[Normality and separability are independent]
\(\mathbb Q(\sqrt[3]2)/\mathbb Q\) is separable because characteristic-zero
fields are perfect, but it is not normal: the two nonreal roots
\[
\omega\sqrt[3]2,\qquad\omega^2\sqrt[3]2
\]
of \(x^3-2\) are absent.  Its normal closure is obtained by adjoining a
primitive cube root of unity \(\omega\).  On the other hand, put
\[
E=\mathbb F_p(t^{1/p}),
\qquad
F=\mathbb F_p(t).
\]
For every \(\alpha\in E\), the Frobenius identity gives
\[
\alpha^p\in F.
\]
Thus the minimal polynomial of \(\alpha\) divides
\[
x^p-\alpha^p=(x-\alpha)^p
\]
in an algebraic closure and has only one distinct root.  Every irreducible
polynomial over \(F\) having a root in \(E\) consequently splits in \(E\) as
a power of a linear factor, so \(E/F\) is normal.  It is not separable,
because the generator \(t^{1/p}\) has derivative-zero minimal polynomial
\(x^p-t\).
\end{example}

The following three familiar extensions summarize the distinctions.  The
middle example has all of its distinct conjugates but not all of them inside
the displayed field; the last loses embeddings because its minimal polynomial
has a repeated root.
\[
\begin{array}{c|c|c|c|c|c}
\text{extension} & \text{degree} & \text{\(F\)-embeddings}
& \text{separable} & \text{normal} & \text{Galois}\\ \hline
\mathbb Q(\sqrt2)/\mathbb Q & 2 & 2 & \text{yes} & \text{yes} & \text{yes}\\
\mathbb Q(\sqrt[3]2)/\mathbb Q & 3 & 3 & \text{yes} & \text{no} & \text{no}\\
\mathbb F_p(t^{1/p})/\mathbb F_p(t) & p & 1 & \text{no} & \text{yes} & \text{no}
\end{array}
\]

\begin{proposition}[Finite Galois criteria]\label{prop:finite-galois-criteria}
A finite extension is Galois precisely when it is the splitting field of a
separable polynomial, equivalently when
\[
|\operatorname{Aut}_F(E)|=[E:F].
\]
\end{proposition}
\begin{proof}
A splitting field is normal by Theorem~\ref{thm:normality-criteria}.
When its defining polynomial is separable, the tower criterion gives
separability.
Conversely, if \(E/F\) is Galois, choose generators
\[
E=F(\alpha_1,\ldots,\alpha_r).
\]
Normality makes every \(m_{\alpha_i,F}\) split in \(E\), and separability
makes each of these polynomials separable.  The product of the distinct
polynomials among them is separable, and its splitting field contains every
generator and is contained in \(E\); hence it is \(E\).

If \(E/F\) is normal and separable, Theorem~\ref{thm:separable-degree} gives
\[
\#\operatorname{Hom}_F(E,\overline F)=[E:F].
\]
Theorem~\ref{thm:normality-criteria} says every such embedding has image
\(E\), so all of them are automorphisms and
\[
|\operatorname{Aut}_F(E)|=[E:F].
\]
Conversely, if this automorphism count holds, then
\[
[E:F]
=|\operatorname{Aut}_F(E)|
\leq\#\operatorname{Hom}_F(E,\overline F)
\leq[E:F].
\]
Equality holds throughout.  Theorem~\ref{thm:separable-degree} gives
separability, and every \(F\)-embedding is already one of the automorphisms.
Its image is therefore \(E\), so Theorem~\ref{thm:normality-criteria} gives
normality.  Thus \(E/F\) is Galois.
\end{proof}

\phantomsection\label{thm:finite-galois-closure}
\begin{applicationtheorem}[Finite Galois closures]
If \(E/F\) is finite separable, then there is a finite Galois extension
\(L/F\) containing \(E\).
\end{applicationtheorem}
\begin{proof}
Choose generators
\[
E=F(\alpha_1,\ldots,\alpha_r),
\]
and let \(f_i\) be the minimal polynomial of \(\alpha_i\) over \(F\).  Each
\(f_i\) is separable because \(E/F\) is separable.  Let \(L\) be the
splitting field over \(F\) of the product of the distinct \(f_i\)'s.  A
splitting field of finitely many polynomials is finite, and \(L/F\) is
normal.  The product is separable, so its splitting field is separable.
Thus \(L/F\) is finite Galois, and it contains every \(\alpha_i\), hence
contains \(E\).
\end{proof}

Such an \(L\) is called a finite Galois closure (or normal Galois closure)
of \(E/F\).  Uniqueness is not needed here.

\section{Finite Fields}
\begin{theorem}[Finite fields]\label{thm:finite-fields}
For each prime power \(q=p^r\), there is, up to isomorphism, one field
\(\mathbb F_q\).  It is the splitting field of
\[
x^q-x
\]
over \(\mathbb F_p\), and \(\mathbb F_q^\times\) is cyclic of order \(q-1\).
Its subfields are exactly \(\mathbb F_{p^d}\) for \(d\mid r\).
\end{theorem}
\begin{proof}
Let \(q=p^r\).  In an algebraic closure of \(\mathbb F_p\), the roots of
\[
x^q-x
\]
form a field: in characteristic \(p\),
\[
(a+b)^q=a^q+b^q,\qquad(ab)^q=a^qb^q,
\]
and nonzero roots remain roots on inversion.  The derivative is \(-1\), so
there are \(q\) distinct roots.  This gives a field with \(q\) elements.

Conversely, if \(K\) has \(q\) elements, then \(K^\times\) has order \(q-1\),
so \(a^{q-1}=1\) for every nonzero \(a\), and \(a^q=a\) for all \(a\in K\).
Thus \(K\) is exactly the root set of \(x^q-x\), proving uniqueness up to
isomorphism and showing that it is a splitting field.  Cyclicity of
\(\mathbb F_q^\times\) is Lemma~\ref{lem:finite-mult-cyclic}.

For subfields, start with an arbitrary field
\[
\mathbb F_p\subseteq K\subseteq\mathbb F_{p^r}
\]
and put \(d=[K:\mathbb F_p]\).  Then \(|K|=p^d\), so every element of \(K\)
satisfies \(x^{p^d}-x=0\).  Thus \(K\) is the unique field with \(p^d\)
elements in the chosen algebraic closure, namely \(\mathbb F_{p^d}\), and
the Tower Law gives \(d\mid r\).  Conversely, if \(d\mid r\), every root of
\[
x^{p^d}-x
\]
lies in \(\mathbb F_{p^r}\), because
\[
x^{p^d}-x\mid x^{p^r}-x
\]
when \(d\mid r\).  To justify the divisibility, write \(r=dk\).  If
\(a^{p^d}=a\), iterating this equality \(k\) times gives
\[
a^{p^r}=a.
\]
Thus every root of \(x^{p^d}-x\) is a root of \(x^{p^r}-x\).  The former
polynomial has derivative \(-1\), hence has no repeated roots, and this
root containment proves the divisibility.  Its \(p^d\) roots form the
required subfield.  This proves both the existence of each listed subfield
and the exhaustion of all subfields.
\end{proof}

\begin{example}[The field with four elements]
The polynomial
\[
x^2+x+1
\]
has no root in \(\mathbb F_2\), so
\[
\mathbb F_4\cong\mathbb F_2[\alpha]/(\alpha^2+\alpha+1).
\]
Its elements are
\[
0,\quad1,\quad\alpha,\quad\alpha+1,
\]
and the relation \(\alpha^2=\alpha+1\) makes multiplication concrete:
\[
(\alpha+1)^2=\alpha^2+1=\alpha.
\]
The Frobenius automorphism fixes \(0,1\) and exchanges the two nontrivial
elements:
\[
\alpha^2=\alpha+1,\qquad(\alpha+1)^2=\alpha.
\]
\end{example}

\begin{example}[The field with eight elements]
By Example~\ref{ex:sep-cubic}, the polynomial \(x^3+x+1\) has no root in
\(\mathbb F_2\) and is irreducible.  Thus
\[
\mathbb F_8\cong\mathbb F_2[\beta]/(\beta^3+\beta+1).
\]
Every element has a unique form
\[
a+b\beta+c\beta^2,
\qquad a,b,c\in\mathbb F_2,
\]
and multiplication is reduced using \(\beta^3=\beta+1\).  For instance,
\[
\beta(\beta^2+1)=\beta^3+\beta=1.
\]
Since \(\beta\ne0,1\) and \(\mathbb F_8^\times\) has prime order \(7\),
Lagrange's theorem implies that \(\beta\) has multiplicative order \(7\).
\end{example}

\begin{proposition}[Frobenius has exact order]\label{prop:frobenius-order}
For \(q=p^r\),
\[
\operatorname{Gal}(\mathbb F_{q^n}/\mathbb F_q)\cong C_n,
\]
generated by Frobenius \(x\mapsto x^q\), whose order is exactly \(n\).
\end{proposition}
\begin{proof}
Frobenius fixes \(\mathbb F_q\), and its \(n\)-th power is the identity on
\(\mathbb F_{q^n}\).  If its \(d\)-th power were the identity for
\(0<d<n\), every element would satisfy \(x^{q^d}=x\), placing the field
among the \(q^d\) roots of \(x^{q^d}-x\), an impossibility since
\(q^n>q^d\).  Thus its \(n\) powers are distinct.  They are all
\(\mathbb F_q\)-automorphisms, and the embedding count bounds the whole
automorphism group by \(n\).
\end{proof}

These cyclic finite Galois groups will become a running example in
Chapter~12.  There the compatible restrictions among the groups
\(\operatorname{Gal}(\mathbb F_{q^n}/\mathbb F_q)\) assemble into the
absolute Galois group of \(\mathbb F_q\).

\begin{example}[Subfields of \(\mathbb F_{p^{12}}\)]
The divisors of \(12\) give all subfields:
\[
\mathbb F_p,\quad
\mathbb F_{p^2},\quad
\mathbb F_{p^3},\quad
\mathbb F_{p^4},\quad
\mathbb F_{p^6},\quad
\mathbb F_{p^{12}}.
\]
For example,
\[
\mathbb F_{p^3}\cap\mathbb F_{p^4}=\mathbb F_p,
\]
since an element in the intersection lies in the unique subfield whose
degree divides both \(3\) and \(4\).
\end{example}

\section{Cyclotomic Extensions}

Cyclotomic fields are the natural meeting point of roots of unity, splitting
fields, and explicit Galois groups.  The multiplicative order of a root of
unity supplies its arithmetic label, while the way automorphisms change that
root will later be read directly from the unit group
\((\mathbb Z/n\mathbb Z)^\times\).

\begin{definition}
A \emph{primitive \(n\)-th root of unity} is an element of multiplicative
order \(n\).  The \(n\)-th cyclotomic polynomial \(\Phi_n(x)\) is the monic
polynomial whose roots are the primitive \(n\)-th roots of unity.
\end{definition}
The finite cyclicity lemma shows that the roots of \(x^n-1\), in
characteristic not dividing \(n\), form a cyclic group.  Partitioning them
by their order gives
\[
x^n-1=\prod_{d\mid n}\Phi_d(x).
\]
Induction in this formula, together with Gauss's Lemma, gives
\(\Phi_n\in\mathbb Z[x]\): after the factors for proper divisors have been
constructed as monic integral polynomials, their product divides the monic
polynomial \(x^n-1\) in \(\mathbb Q[x]\), and polynomial division followed by
Gauss's Lemma makes the monic quotient integral.  The number of primitive
roots gives
\[
\deg\Phi_n=\varphi(n).
\]
For example,
\[
\begin{aligned}
\Phi_1&=x-1, & \Phi_2&=x+1, & \Phi_3&=x^2+x+1,\\
\Phi_4&=x^2+1, & \Phi_5&=x^4+x^3+x^2+x+1,\\
\Phi_6&=x^2-x+1, & \Phi_8&=x^4+1.
\end{aligned}
\]

\begin{theorem}[Cyclotomic irreducibility]
\(\Phi_n(x)\) is irreducible over \(\mathbb Q\).
\end{theorem}
\begin{proof}
Let \(\zeta\) be primitive and let \(g\in\mathbb Z[x]\) be its monic
irreducible polynomial.  Write
\[
\Phi_n=gh
\]
in \(\mathbb Z[x]\).  Let \(p\nmid n\) be prime.  Suppose that \(\zeta^p\)
is not a root of \(g\).  It is a primitive root and hence a root of \(h\),
so \(h(\zeta^p)=0\).  Therefore \(\zeta\) is a root of \(h(x^p)\), and
minimality of \(g\) gives
\[
g(x)\mid h(x^p)
\]
in \(\mathbb Q[x]\), hence in \(\mathbb Z[x]\) by Gauss's Lemma.  Reducing
modulo \(p\), Frobenius gives
\[
\overline{h(x^p)}=\overline{h(x)}^{\,p}.
\]
Thus an irreducible factor of \(\bar g\) also divides \(\bar h\), so
\(\overline{\Phi_n}=\bar g\bar h\) has a repeated irreducible factor.
But \(\Phi_n\) divides \(x^n-1\), whereas
\[
(x^n-1)'=nx^{n-1}
\]
is relatively prime to \(x^n-1\) modulo \(p\).  This contradicts the
derivative criterion.  Hence \(\zeta^p\) is a root of \(g\) for every prime
\(p\nmid n\).  Let \(a>0\) with \((a,n)=1\), and factor
\[
a=p_1\cdots p_s.
\]
No \(p_i\) divides \(n\).  Repeatedly applying the preceding conclusion
shows that \(\zeta^a\) is a root of \(g\).  Thus every primitive root is a
root of \(g\), and consequently \(g=\Phi_n\).
\end{proof}

If \(\zeta_n\) is primitive, irreducibility gives
\[
[\mathbb Q(\zeta_n):\mathbb Q]=\varphi(n).
\]
It is the splitting field of \(x^n-1\), hence is finite Galois.  The map
\[
(\mathbb Z/n\mathbb Z)^\times\longrightarrow
\operatorname{Gal}(\mathbb Q(\zeta_n)/\mathbb Q),
\qquad a\longmapsto(\zeta_n\mapsto\zeta_n^a)
\]
is an isomorphism.  Indeed, \(\zeta_n^a\) is primitive exactly when
\((a,n)=1\), and then it is a root of the minimal polynomial \(\Phi_n\).
The root-extension lemma therefore gives a unique \(\mathbb Q\)-embedding
sending \(\zeta_n\) to \(\zeta_n^a\); finite dimensionality makes it an
automorphism.  Composition corresponds to multiplication of exponents, and
the map is injective because an automorphism is determined by \(\zeta_n\).
Conversely every automorphism sends \(\zeta_n\) to a primitive root, so it
has this form.  Roots of unity will reappear in the later discussion of
radicals.

\begin{example}[Small cyclotomic Galois groups]
For \(n=3\) and \(n=4\),
\[
(\mathbb Z/n\mathbb Z)^\times\cong C_2,
\]
so both \(\mathbb Q(\zeta_3)\) and
\(\mathbb Q(\zeta_4)=\mathbb Q(i)\) are quadratic Galois extensions.  For
\(n=5\),
\[
(\mathbb Z/5\mathbb Z)^\times\cong C_4,
\]
whereas
\[
(\mathbb Z/8\mathbb Z)^\times\cong C_2\oplus C_2.
\]
Thus the degree-four fields \(\mathbb Q(\zeta_5)\) and
\(\mathbb Q(\zeta_8)\) have nonisomorphic Galois groups.
\end{example}

\section*{Exercises}
\subsection*{Basic}
\begin{exercise}[Basic] Determine whether \(x^4-2x^2+1\) is separable over
\(\mathbb Q\), and explain the answer using the derivative criterion.
\end{exercise}
\begin{exercise}[Basic] Use the derivative criterion to decide whether
\[
x^p-x\in\mathbb F_p[x]
\]
is separable.
\end{exercise}
\begin{exercise}[Basic] Show that a finite field is perfect directly from
the fact that its Frobenius map is bijective.
\end{exercise}

\subsection*{Core}
\begin{exercise}[Core] Test normality and separability of
\(\mathbb Q(\sqrt[3]2)/\mathbb Q\).
\end{exercise}
\begin{exercise}[Core] Let \(E/F\) be finite.  Explain why
\[
[E:F]_s=[E:F]
\]
when \(E/F\) is generated by a tower of separable simple extensions.
\end{exercise}
\begin{exercise}[Core] Find a primitive element for
\[
\mathbb Q(\sqrt2,\sqrt5)/\mathbb Q.
\]
\end{exercise}
\begin{exercise}[Core] Verify that \(x^3+x+1\) is irreducible over
\(\mathbb F_2\), and use it to list the elements of \(\mathbb F_8\).
\end{exercise}
\begin{exercise}[Core] Determine the subfields of \(\mathbb F_{p^{12}}\).
\end{exercise}
\begin{exercise}[Core] Prove that every finite extension of a finite field
is Galois.
\end{exercise}
\begin{exercise}[Core] Compute the Frobenius orbit of a nonzero element of
\(\mathbb F_8\) not in \(\mathbb F_2\).
\end{exercise}

\subsection*{Further}
\begin{exercise}[Further] Show that
\[
\Phi_{10}(x)=x^4-x^3+x^2-x+1.
\]
\end{exercise}
\begin{exercise}[Further] Show that \(\mathbb Q(\zeta_5)\) has Galois group
isomorphic to \(C_4\).
\end{exercise}
\begin{exercise}[Further] Determine
\(\operatorname{Gal}(\mathbb Q(\zeta_8)/\mathbb Q)\).
\end{exercise}
\begin{exercise}[Further] Prove that the splitting field of
\[
x^4-2
\]
over \(\mathbb Q\) is normal and separable, and determine its degree.
\end{exercise}

\subsection*{Looking Ahead}
\begin{exercise}[Challenging] Let \(E/F\) be finite and separable.  Use the
embedding-count theorem to show that \(E/F\) is normal exactly when every
\(F\)-embedding \(E\to\overline F\) is an automorphism of \(E\).
\end{exercise}
\begin{exercise}[Looking Ahead] If \(E/F\) is finite Galois, show that
\[
F\subseteq E^{\operatorname{Aut}_F(E)}.
\]
The reverse inclusion is the starting point of the Galois correspondence.
\end{exercise}
